% ========================================================================
% فصل ۱۴: جمع‌بندی
% خروجی عامل sam-master: sam-work/sections/14-conclusion.tex
% منابع مجاز: hozouriOverviewVANETVehicular2023، xuInternetVehiclesBig2018،
%   farsimadanReviewSecurityChallenges2025، rahmawatiagustinaSystematicLiteratureReview2025،
%   razaBlockchainBasedReputationTrust2024، ghovanlooyghajarSBTMSScalableBlockchain2021،
%   IOTATANGLE20، sealeyIOTATangle202022، silvanoIotaTangleCryptocurrency2020،
%   mazzocca2025survey، feraudoDIVADIDbasedReputation2024،
%   al-aniPrivacySafetyImprovement2023، luPseudonymChangingSocial2012،
%   fengBPASBlockchainAssistedPrivacyPreserving2020، mariaBAIVEfficientBlockchainBased2022،
%   linGSISSecurePrivacyPreserving2007، SecuringVehicularAd
% ========================================================================

\section{جمع‌بندی}

در این گزارش، چالش‌ها و راهکارهای امنیتی در شبکه‌ی اقتضایی خودرویی (\lr{VANET}\LTRfootnote{\lr{Vehicular Ad hoc Network (VANET)}}) به‌طور جامع مورد بررسی قرار گرفت. شبکه‌ی اقتضایی خودرویی به‌عنوان زیرساخت ارتباطی نوین در سامانه‌های حمل‌ونقل هوشمند (\lr{ITS}\LTRfootnote{\lr{Intelligent Transportation System (ITS)}})، نقشی حیاتی در بهبود ایمنی ترافیک و کاهش تصادفات ایفا می‌کند. این فصل یافته‌های فصل‌های پیشین را بازنویسی نمی‌کند، بلکه میان آن‌ها پیوند برقرار می‌کند: آسیب‌پذیری‌های ذاتی و حملات امنیتی، رمزنگاری و زیرساخت کلید عمومی، حریم خصوصی و مدیریت اعتماد، زنجیره‌ی بلوکی و فناوری‌های دفتر کل توزیع‌شده، هویت غیرمتمرکز و اعتبارنامه‌های تأییدپذیر، تعادل حریم خصوصی و ایمنی و چشم‌انداز آینده در هفت زیربخش جمع‌بندی می‌شوند.

\subsection{آسیب‌پذیری‌های ذاتی و حملات امنیتی}

همان‌طور که در مقدمه اشاره شد، ریشه‌ی بسیاری از این چالش‌ها در ویژگی‌های ساختاری شبکه است؛ بنابراین نخستین زیربخش به آسیب‌پذیری‌های ذاتی و حملات امنیتی می‌پردازد. مرورهای کلی، تحرک بالا، توپولوژی پویا، اتصال پراکنده و تکه‌تکه‌شدن مکرر شبکه را از ویژگی‌های ذاتی شبکه‌ی اقتضایی خودرویی می‌دانند \cite{hozouriOverviewVANETVehicular2023, xuInternetVehiclesBig2018}. پیام‌های ایمنی دوره‌ای که موقعیت، سرعت و جهت خودرو را حمل می‌کنند، به‌صورت متن ساده و با نرخ ۱ تا ۱۰ هرتز پخش می‌شوند \cite{al-aniPrivacySafetyImprovement2023} و ماهیت پخش‌کننده‌ی این پیام‌ها آن‌ها را در معرض بازپخش، جعل و استراق سمع قرار می‌دهد \cite{farsimadanReviewSecurityChallenges2025}. مرور فرسیمادان و همکاران حملات را در پنج دسته بررسی می‌کنند \cite{farsimadanReviewSecurityChallenges2025} و رایا و هوبو مهاجمان را در چهار بعد طبقه‌بندی می‌کنند: داخلی در برابر خارجی، خبیث در برابر عقلانی، فعال در برابر منفعل و محلی در برابر گسترده \cite{SecuringVehicularAd}. شدت این تهدیدات از آنجاست که مهاجمی که کنترل خودروها را به دست آورد، می‌تواند تصادفاتی ایجاد کند که جان رانندگان را به خطر بیندازد \cite{farsimadanReviewSecurityChallenges2025} و حمله‌ی سیبل\LTRfootnote{\lr{Sybil attack}} از جمله تهدیداتی است که طرح‌های شبه‌نام باید در برابر آن مقاومت کنند \cite{rahmawatiagustinaSystematicLiteratureReview2025}. نتیجه‌گیری کلیدی این زیربخش آن است که آسیب‌پذیری، عارضه‌ی طراحی نیست بلکه نتیجه‌ی ناگزیر باز بودن شبکه و الزامات ایمنی است و از این‌رو دفاع باید لایه‌بندی باشد.

\subsection{رمزنگاری و زیرساخت کلید عمومی}

آسیب‌پذیری‌های فوق ریشه در ساختار باز و پویای شبکه دارند؛ نخستین لایه‌ی دفاع در برابر آن‌ها، رمزنگاری و زیرساخت کلید عمومی (\lr{PKI}\LTRfootnote{\lr{Public Key Infrastructure (PKI)}}) است. در مدل استاندارد، هر پیام با امضای دیجیتال و گواهی صادرشده توسط مرجع صدور گواهی (\lr{CA}\LTRfootnote{\lr{Certificate Authority (CA)}}) تأیید می‌شود \cite{SecuringVehicularAd}. امضای دیجیتال با وجود سرباری نزدیک ۱۴۰ بایت، «راحت‌ترین و قابل‌اعتمادترین» روش تأیید پیام تلقی می‌شود و کلید ۲۲۴ بیتی \lr{ECC}\LTRfootnote{\lr{Elliptic Curve Cryptography (ECC)}} تقریباً معادل کلید ۲۰۴۸ بیتی \lr{RSA}\LTRfootnote{\lr{RSA}} است \cite{SecuringVehicularAd}. طرح‌های تازه‌تر تلاش می‌کنند سربار و وابستگی به مرجع صادرکننده را کاهش دهند: پروتکل \lr{GSIS}\LTRfootnote{\lr{GSIS}} امضای گروهی و امضای مبتنی بر شناسه را ترکیب می‌کند و امکان ردیابی شرطی را فراهم می‌کند \cite{linGSISSecurePrivacyPreserving2007}؛ طرح \lr{BAIV}\LTRfootnote{\lr{BAIV}} هزینه‌ی تصدیق یک کاربر و یک واحد کنارجاده‌ای (\lr{RSU}\LTRfootnote{\lr{RSU}}) را به ۱۲٫۵۴ میلی‌ثانیه، یعنی حدود ۸۰ تصدیق در ثانیه، می‌رساند \cite{mariaBAIVEfficientBlockchainBased2022}؛ و طرح \lr{BPAS}\LTRfootnote{\lr{BPAS}} مرجع ثبت‌نام آنلاین را جز در آغاز سامانه و ثبت‌نام خودرو حذف می‌کند، ردیابی شرطی و لغو پویا را فراهم می‌کند و بر پلتفرم \lr{Hyperledger Fabric}\LTRfootnote{\lr{Hyperledger Fabric}} ارزیابی شده است \cite{fengBPASBlockchainAssistedPrivacyPreserving2020}. با این حال، طرح \lr{BAIV} همچنان به مرجع قابل اعتماد (\lr{TA}\LTRfootnote{\lr{Trusted Authority (TA)}}) برای ثبت‌نام آفلاین، نگاشت هویت و لغو نیاز دارد \cite{mariaBAIVEfficientBlockchainBased2022} و لغو اعتبارنامه‌ها «شاید دشوارترین جنبه‌ی امنیت شبکه‌ی اقتضایی خودرویی» است؛ برای رانندگی دو ساعت در روز، حدود ۴۳٬۸۰۰ کلید در سال و حدود ۴٫۲ مگابایت فهرست لغو تولید می‌شود \cite{SecuringVehicularAd} و مدیریت کلید را به‌دلیل تغییر توپولوژی «بسیار دشوار» می‌کند \cite{hozouriOverviewVANETVehicular2023}. محدودیت دیگر آن است که سربار محاسباتی، برخی راهکارهای امنیتی را بر ظرفیت محدود محاسباتی خودروها غیرقابل اجرا می‌کند \cite{farsimadanReviewSecurityChallenges2025} و محرمانگی برای پیام‌های ایمنی عمومی \lr{V2X}\LTRfootnote{\lr{Vehicle-to-Everything (V2X)}} اولویت نخست نیست \cite{farsimadanReviewSecurityChallenges2025, SecuringVehicularAd}. در بلندمدت نیز الگوریتم‌های \lr{RSA} و \lr{ECC} در برابر تهدید کوانتومی آسیب‌پذیرند و امضای یک‌بارمصرف مقاوم در برابر کوانتوم، جایگزینی را مطرح کرده است \cite{sealeyIOTATangle202022}. مهم‌تر از همه، رمزنگاری تهدیدات بیرونی را محور قرار می‌دهد و در برابر حملات داخلی توسط خودروهای تأییدشده آسیب‌پذیر است \cite{farsimadanReviewSecurityChallenges2025}، زیرا مهاجم داخلی کلید عمومی تأییدشده‌ی در اختیار دارد \cite{SecuringVehicularAd}؛ بنابراین رمزنگاری لایه‌ای لازم اما ناکافی است.

\subsection{حریم خصوصی و مدیریت اعتماد}

رمزنگاری، یکپارچگی و اصالت پیام را تأمین می‌کند، اما هویت فرستنده را در معرض افشا قرار می‌دهد؛ لایه‌ی بعدی، حریم خصوصی و مدیریت اعتماد، این خلأ را پر می‌کند. شبه‌نام\LTRfootnote{\lr{Pseudonym}} شناسه‌ای موقت است که هویت واقعی خودرو را پنهان می‌کند \cite{rahmawatiagustinaSystematicLiteratureReview2025} و حریم خصوصی هویت، پیش‌شرط حریم خصوصی موقعیت است \cite{luPseudonymChangingSocial2012}. مؤثر بودن این مکانیزم به شرط حساسی مکان و زمان بستگی دارد: تغییر شبه‌نام تنها در زمان و مکان مناسب کارآمد است و در غیر آن بی‌فایده می‌شود \cite{luPseudonymChangingSocial2012} و بازه‌های ثابت و قابل پیش‌بینی، حمله‌ی پیوند\LTRfootnote{\lr{Linkability attack}} را تسهیل می‌کند \cite{rahmawatiagustinaSystematicLiteratureReview2025}. چرخه‌ی حیات شبه‌نام پنج مرحله‌ی صدور، استفاده، تغییر، حل و لغو دارد \cite{rahmawatiagustinaSystematicLiteratureReview2025}، اما در مروری بر ۱۱۳ مطالعه‌ی ۲۰۱۹ تا ۲۰۲۴، تنها ۱۴ مطالعه، یعنی حدود ۱۲ درصد، همه‌ی این مراحل را پیاده کرده بودند \cite{rahmawatiagustinaSystematicLiteratureReview2025}. گزارش فنی \lr{ETSI TR 103 415}\LTRfootnote{\lr{European Telecommunications Standards Institute (ETSI) Technical Report 103 415}} شش راهبرد برای تغییر شبه‌نام برشمرده است \cite{rahmawatiagustinaSystematicLiteratureReview2025}. یکی از رویکردهای عملی، تغییر شبه‌نام در نقاط اجتماعی است: چون خودروها در نقاطی مانند تقاطع هنگام چراغ قرمز یا پارکینگ رایگان به‌طور موقت جمع می‌شوند، این نقاط به‌طور طبیعی به ناحیه‌ی اختلاط تبدیل می‌شوند \cite{luPseudonymChangingSocial2012}. رویکرد دیگر، طرح \lr{SRPS}\LTRfootnote{\lr{Safety-related Privacy Scheme (SRPS)}} است که با ردیاب چندهدفه (\lr{MTT}\LTRfootnote{\lr{Multi-Target Tracking (MTT)}}) مبتنی بر فیلتر کالمن\LTRfootnote{\lr{Kalman filter}}، انتساب احتمالاتی داده‌ی نزدیک‌ترین همسایه (\lr{NNPDA}\LTRfootnote{\lr{Nearest Neighbour Probabilistic Data Association (NNPDA)}}) و نگهداشت مسیر ردیابی، وضعیت همسایگان را حتی در دوره‌ی سکوت پیش‌بینی می‌کند، در صورت پیش‌بینی تصادف از سکوت خارج می‌شود و در صورت احتمال اختلاط زمینه، بدون ورود به سکوت شبه‌نام عوض می‌کند \cite{al-aniPrivacySafetyImprovement2023}. در کنار این مکانیزم‌ها، اعتماد\LTRfootnote{\lr{Trust}} به‌عنوان معیار امنیتی تکمیلی، اطمینان یک نهاد به نهاد دیگر است \cite{farsimadanReviewSecurityChallenges2025} و روش جمع وزنی شاخص‌های اعتماد، رایج‌ترین راهبرد مدیریت اعتماد است \cite{farsimadanReviewSecurityChallenges2025}. اما گره‌های خبیث پیشرفته با جابه‌جایی میان رفتار عادی و خبیث، شناسایی را دشوار می‌کنند \cite{farsimadanReviewSecurityChallenges2025} و امتیازات اعتماد در زنجیره‌ی بلوکی نگهداری و میان واحدهای کنارجاده‌ای به اشتراک گذاشته می‌شوند \cite{razaBlockchainBasedReputationTrust2024}. جمع‌بندی این زیربخش آن است که شبه‌نام و اعتماد مکمل‌اند: نخستی هویت را پنهان می‌کند و دومی رفتار را ارزیابی می‌کند، اما هر دو در عمل به مدیریت متمرکز وابسته‌اند.

\subsection{زنجیره بلوکی و فناوری‌های دفتر کل توزیع‌شده}

صدور شبه‌نام و مدیریت اعتماد در طرح‌های موجود عمدتاً به مرجع قابل اعتماد وابسته‌اند و این وابستگی نقطه شکست واحد متمرکز ایجاد می‌کند؛ زنجیره‌ی بلوکی و فناوری‌های دفتر کل توزیع‌شده به‌عنوان راهکار غیرمتمرکز برای این مسئله مطرح شده‌اند. زنجیره‌ی بلوکی شبکه‌ی همتا به همتا‌ای برای تراکنش امن بدون طرف سوم قابل اعتماد است و بلوک‌ها با درج درهم‌ساز بلوک قبل به‌صورت پشت‌سرهم پیوند می‌خورند \cite{razaBlockchainBasedReputationTrust2024}؛ مزایای آن غیرمتمرکزسازی با حذف نقطه شکست واحد، تغییرناپذیری و قابلیت ردیابی است \cite{razaBlockchainBasedReputationTrust2024}. با این حال، اثبات کار (\lr{PoW}\LTRfootnote{\lr{Proof of Work (PoW)}}) محاسباتی سنگین، کند و با تأخیر بالاست و برای شبکه‌های پویایی مانند شبکه‌ی اقتضایی خودرویی نامناسب است و پلتفرم‌هایی مانند اتریوم و هایپرلدجر نیز ظرفیت محدودی دارند \cite{razaBlockchainBasedReputationTrust2024}؛ برآوردی که در یکی از مقالات ذکر شده، حدود ۷ تراکنش در ثانیه برای بیت‌کوین و ۱۵ تراکنش در ثانیه برای اتریوم است، هرچند منبع اولیه‌ی این ارقام در آن مقاله ذکر نشده \cite{IOTATANGLE20}. شاردینگ\LTRfootnote{\lr{Sharding}} دفتر کل را به قطعه‌هایی تقسیم می‌کند و وظایف استخراج را میان کمیته‌ها توزیع می‌کند \cite{ghovanlooyghajarSBTMSScalableBlockchain2021}. در سامانه‌ی \lr{SBTMS}\LTRfootnote{\lr{SBTMS}}، قابلیت اطمینان پیام با استنباط بیزی ارزیابی می‌شود، تحمل خطای بیزانسی عملی (\lr{PBFT}\LTRfootnote{\lr{Practical Byzantine Fault Tolerance (PBFT)}}) در هر کمیته و سپس در کمیته‌ی نهایی اجرا می‌شود و اثبات کار تنها برای هویت گذرای واحدهای کنارجاده‌ای نگه داشته می‌شود \cite{ghovanlooyghajarSBTMSScalableBlockchain2021}؛ زمان تولید بلوک نیز با افزایش تعداد خودروها به‌طور نزدیک به خطی رشد می‌کند و بر اساس خوانش نمودار مقاله، برای ۳۰۰۰ خودرو این زمان در اثبات کار حدود ۱۴ ثانیه و در شاردینگ حدود ۲ ثانیه است \cite{ghovanlooyghajarSBTMSScalableBlockchain2021}. تنگل \lr{IOTA}\LTRfootnote{\lr{IOTA}} دفتر کل توزیع‌شده‌ای مبتنی بر گراف جهت‌دار غیرمدور (\lr{DAG}\LTRfootnote{\lr{Directed Acyclic Graph (DAG)}}) برای اینترنت اشیا است \cite{silvanoIotaTangleCryptocurrency2020, sealeyIOTATangle202022}؛ در نسخه‌های سنتی هر تراکنش جدید دو تراکنش قبلی را تأیید می‌کند \cite{silvanoIotaTangleCryptocurrency2020} و در نسخه‌ی ۲٫۰ به ۲ تا ۸ پیام ارجاع می‌دهد \cite{IOTATANGLE20, sealeyIOTATangle202022}. این فناوری بدون کارمزد و بدون نیاز به استخراج‌کننده است \cite{sealeyIOTATangle202022, IOTATANGLE20}؛ هماهنگ‌کننده‌ای که در نسخه‌های سنتی وجود داشت \cite{silvanoIotaTangleCryptocurrency2020} در نسخه‌ی ۲٫۰ حذف شده است \cite{IOTATANGLE20, sealeyIOTATangle202022}، اجماع احتمالاتی سریع (\lr{FPC}\LTRfootnote{\lr{Fast Probabilistic Consensus (FPC)}}) بی‌رهبر است \cite{sealeyIOTATangle202022} و سامانه‌ی شهرت مانا (\lr{Mana}\LTRfootnote{\lr{Mana}}) در برابر حمله‌ی سیبل و کسوف به کار می‌رود \cite{sealeyIOTATangle202022}. پیاده‌سازی آزمایشی \lr{Nectar}\LTRfootnote{\lr{Nectar}} تا ۱۰۰۰ تراکنش در ثانیه با زمان تأیید میانگین ۱۰ تا ۱۲ ثانیه را پشتیبانی می‌کند، اما این ارقام از منابع بنیاد \lr{IOTA} نقل شده و ارزیابی مستقل نیست \cite{sealeyIOTATangle202022} و توسعه‌ی نسخه‌ی ۲٫۰ همچنان ادامه دارد \cite{IOTATANGLE20}؛ نسخه‌های پیش از آن نیز «نیمه‌غیرمتمرکز» توصیف شده‌اند \cite{silvanoIotaTangleCryptocurrency2020}. سامانه‌ی \lr{DIVA}\LTRfootnote{\lr{DIVA}} نیز از دفتر کل \lr{DAG} با دسترسی محدود استفاده می‌کند \cite{feraudoDIVADIDbasedReputation2024}. نکته‌ی تقاطعی این زیربخش آن است که در یک مقایسه، هماهنگ‌کننده‌ی تنگل نقطه شکست واحد تلقی شده و شاردینگ به‌عنوان جایگزینی کاملاً توزیع‌شده معرفی شده است \cite{ghovanlooyghajarSBTMSScalableBlockchain2021}. در مجموع، هیچ مکانیزم اجماعی هنوز همه‌ی الزامات بلادرنگ شبکه‌ی اقتضایی خودرویی را برآورده نکرده است: زنجیره‌ی بلوکی می‌تواند پیچیدگی زمانی بیشتر ایجاد کند \cite{farsimadanReviewSecurityChallenges2025} و اجماع‌های جایگزین مانند اثبات حصه (\lr{PoS}\LTRfootnote{\lr{Proof of Stake (PoS)}}) و \lr{PBFT} می‌توانند تأخیر را کاهش دهند اما ممکن است نیازهای سخت‌گیرانه‌ی تأخیر کم در کاربردهای بلادرنگ این شبکه را برآورده نکنند \cite{razaBlockchainBasedReputationTrust2024}.

\subsection{هویت غیرمتمرکز و اعتبارنامه‌های تأییدپذیر}

دفتر کل توزیع‌شده، زیرساخت ذخیره‌سازی و اجماع را فراهم می‌کند؛ لایه‌ی هویت بر آن، هویت غیرمتمرکز و اعتبارنامه‌های تأییدپذیر را می‌سازد. هویت خودمختار (\lr{SSI}\LTRfootnote{\lr{Self-Sovereign Identity (SSI)}}) که در ۲۰۱۲ معرفی شد، بر دفترهای کل توزیع‌شده، شناسه‌های غیرمتمرکز و اعتبارنامه‌های تأییدپذیر استوار است \cite{mazzocca2025survey}. شناسه‌ی غیرمتمرکز (\lr{DID}\LTRfootnote{\lr{Decentralized Identifier (DID)}}) طرح شناسه‌ای جهانی و یکتا بر پایه‌ی رمزنگاری است و استاندارد کنسرسیوم وب جهان‌گستر (\lr{W3C}\LTRfootnote{\lr{World Wide Web Consortium (W3C)}}) است \cite{mazzocca2025survey} و اعتبارنامه‌ی تأییدپذیر (\lr{VC}\LTRfootnote{\lr{Verifiable Credential (VC)}}) مشخصات \lr{W3C} برای ادعاهای رمزنگاری قابل تأیید و ضد جعل است \cite{mazzocca2025survey}. این چارچوب نقش‌های دارنده\LTRfootnote{\lr{Holder}}، صادرکننده\LTRfootnote{\lr{Issuer}} و تأییدکننده\LTRfootnote{\lr{Verifier}} و ثبت‌کننده‌ی داده تأییدپذیر (\lr{VDR}\LTRfootnote{\lr{Verifiable Data Registry (VDR)}}) را تعریف می‌کند \cite{mazzocca2025survey} و افشای انتخابی\LTRfootnote{\lr{Selective disclosure}} را از طریق ادعاهای تکی، مقادیر درهم‌ساز، اثبات دانش‌صفر (\lr{ZKP}\LTRfootnote{\lr{Zero-Knowledge Proof (ZKP)}}) و امضاهای افشای انتخابی انجام می‌دهد؛ اعتبارنامه‌ی \lr{SD-JWT}\LTRfootnote{\lr{Selective Disclosure JSON Web Token (SD-JWT)}} راه‌حل پیشروی کنونی است \cite{mazzocca2025survey}. در کاربرد خودرویی، چارچوب ارتباطات خودرویی غیرمتمرکز (\lr{D-V2X}\LTRfootnote{\lr{Decentralized V2X (D-V2X)}}) با \lr{DID} و \lr{VC} یکپارچگی، تأیید هویت، محرمانگی و حریم خصوصی را بدون نهاد متمرکز قابل اعتماد ممکن می‌سازد: خودرو چندین شناسه‌ی غیرمتمرکز ثبت می‌کند که با اعتبارنامه‌های مبتنی بر اثبات دانش‌صفر به شناسه‌ی اصلی پیوند می‌خورند و با ساختن دو ارائه تأییدپذیر (\lr{VP}\LTRfootnote{\lr{Verifiable Presentation (VP)}}) شبه‌نام ایجاد می‌کند \cite{mazzocca2025survey}. این مدل مقررات حفاظت عمومی از داده‌ها (\lr{GDPR}\LTRfootnote{\lr{General Data Protection Regulation (GDPR)}}) را پشتیبانی می‌کند، زیرا کنترل کامل داده را به فرد می‌دهد بدون آنکه اعزام آن به نهاد متمرکز را اجباری کند \cite{mazzocca2025survey} و با مقررات اتحادیه‌ی اروپا با شماره‌ی ۲۰۲۴/۱۱۸۳ که در مه ۲۰۲۴ چارچوب شناسه‌ی دیجیتال اروپایی را تأسیس کرد، هم‌سو است \cite{mazzocca2025survey}. سامانه‌ی \lr{DIVA} نمونه‌ای عملی است: مرجع قابل اعتماد، شناسه‌ی غیرمتمرکز و سند شناسه\LTRfootnote{\lr{DID Document}} را بر دفتر کل ثبت می‌کند و اعتبارنامه صادر می‌کند، خودرو به هر پیام یک ارائه تأییدپذیر پیوست می‌کند، امتیاز شهرت با تراکنش بدون ارزش\LTRfootnote{\lr{Zero-value Transaction}} نگهداری و روی گره‌های لبه محاسبه می‌شود \cite{feraudoDIVADIDbasedReputation2024}. دقت تشخیص منابع خبیث حدود ۹۴ درصد با ۳۰ درصد و حدود ۸۹ درصد با ۴۰ درصد منابع خبیث (آستانه‌ی میانگین با بیتای بزرگ‌تر از ۰٫۳) بوده و نرخ تشخیص مثبت (\lr{TPR}\LTRfootnote{\lr{True Positive Rate (TPR)}}) ۹۹٫۹۳ درصد و نرخ تشخیص منفی (\lr{TNR}\LTRfootnote{\lr{True Negative Rate (TNR)}}) ۷۸٫۱۴ درصد (آستانه‌ی مد) گزارش شده است \cite{feraudoDIVADIDbasedReputation2024}. با این حال، سامانه‌ی \lr{DIVA} همچنان به مرجع قابل اعتماد نیاز دارد، دفتر کل آن با دسترسی محدود است، شناسه‌های غیرمتمرکز در آن «قابل نگاشت به هویت کاربر نیستند» (یعنی سازوکار افشای هویت واقعی ندارد)، ارزیابی تنها در یک سناریو و فقط بر پیام‌های رویدادمحور انجام شده و تأخیر فقط برای ارتباط یک‌به‌یک سنجیده شده است \cite{feraudoDIVADIDbasedReputation2024}. چالش‌های کلان این حوزه نیز استانداردسازی (بدون توافق درباره‌ی الگوریتم‌ها و مدیریت کلید)، مقیاس‌پذیری، پاسخ‌گویی (از جمله احراز هویت مشتری و مبارزه با پولشویی) و لغو است؛ برای لغو تنها یک مشخصه‌ی پیشنهادی وجود دارد که هنوز استاندارد \lr{W3C} نیست و در محیط‌های پرترافیک و برای دستگاه‌های محدود کمتر بررسی شده است \cite{mazzocca2025survey}. نتیجه‌گیری این زیربخش آن است که این فناوری‌ها موعود هستند اما هنوز به مراجع قابل اعتماد و دفترهای کل مجوزدار وابسته‌اند.

\subsection{تعادل حریم خصوصی و ایمنی}

هویت غیرمتمرکز، کنترل افشای اطلاعات را به کاربر می‌دهد؛ با این حال، این کنترل باید با الزامات ایمنی تعادل شود. پیام‌های ایمنی دوره‌ای متن‌ساده، امکان ردیابی مکانی‌زمانی خودروها را فراهم می‌کنند \cite{al-aniPrivacySafetyImprovement2023}. دوره‌های سکوت، حریم خصوصی را بهتر می‌کنند اما «تصادف بالقوه‌ای در این دوره قابل پیشگیری نیست» \cite{al-aniPrivacySafetyImprovement2023} و تغییر بیشتر شبه‌نام و سکوت طولانی‌تر، سربار امنیتی بیشتر و از دست رفتن پیام‌ها را به دنبال دارد \cite{al-aniPrivacySafetyImprovement2023}. در ارزیابی طرح \lr{SRPS}، ۹٫۸۱، ۹٫۸۷ و ۹٫۹۲ پیام ایمنی در ثانیه به دست آمد (در برابر ۱۰٫۰۰ برای طرح \lr{PPC}\LTRfootnote{\lr{PPC}}، ۹٫۳۳ تا ۹٫۴۱ برای \lr{CAPS}\LTRfootnote{\lr{CAPS}}، ۶٫۱۹ تا ۶٫۴۷ برای \lr{SLOW}\LTRfootnote{\lr{SLOW}} و ۷٫۲۲ تا ۷٫۶۵ برای \lr{RSP}\LTRfootnote{\lr{RSP}})؛ درصد خودروهای قابل ردیابی به ۴۰، ۳۱ و ۲۰ درصد رسید (در برابر ۷۱، ۶۳ و ۵۲ درصد برای \lr{CAPS} و تا ۹۴ درصد برای \lr{PPC})؛ سردرگمی مهاجم به ۶۹، ۷۷ و ۸۴ درصد افزایش یافت (در برابر ۲۶، ۳۸ و ۴۱ درصد برای \lr{CAPS}) و ۴۳، ۶۸ و ۹۳ تصادف بالقوه اجتناب‌شده (\lr{PAA}\LTRfootnote{\lr{Potentially Avoided Accidents (PAA)}}) گزارش شد \cite{al-aniPrivacySafetyImprovement2023}. نویسندگان اعلام می‌کنند طرح \lr{SRPS} «بهترین تعادل میان سه مسئله‌ی کلیدی حریم خصوصی، ایمنی و کارایی» را به دست آورده است \cite{al-aniPrivacySafetyImprovement2023}، هرچند این مقادیر از روی برچسب نمودارها خوانده شده‌اند و \lr{PAA} تعداد تصادف‌های مورد انتظار است نه تصادف‌های واقعاً شبیه‌سازی‌شده \cite{al-aniPrivacySafetyImprovement2023}. کشمکش بنیادی، کشمکش گمنامی و پاسخ‌گویی است: مرجع ردیابی (\lr{TRA}\LTRfootnote{\lr{Trace Authority (TRA)}}) تعادل میان حفاظت از حریم خصوصی و پاسخ‌گویی برقرار می‌کند \cite{rahmawatiagustinaSystematicLiteratureReview2025} و در حریم خصوصی شرطی، مرجع مجاز در رویدادهای اختلافی هویت فرستنده را افشا می‌کند \cite{linGSISSecurePrivacyPreserving2007, luPseudonymChangingSocial2012}. نمونه‌ی وارونه نیز هشداردهنده است: طرحی که همه‌ی خودروها را همزمان وارد دوره‌ی سکوت می‌کند، سردرگمی ۱۰۰ درصد دارد اما «شبکه‌ی اقتضایی خودرویی با ورود همه‌ی خودروها به دوره‌ی سکوت به‌طور همزمان از کار می‌افتد» \cite{al-aniPrivacySafetyImprovement2023}. پس تعادل حریم خصوصی و ایمنی نقطه‌ی ثابتی نیست، بلکه تعادلی زمینه‌ای و پویاست.

\subsection{چشم‌انداز آینده}

تعادل حریم خصوصی و ایمنی، پیامدی از شکاف‌های پژوهشی کنونی است؛ چشم‌انداز آینده این شکاف‌ها را بررسی می‌کند. نخست، تکمیل چرخه‌ی حیات شبه‌نام بدون افت عملکرد و با کمینه‌سازی تأخیر، چالشی باز است \cite{rahmawatiagustinaSystematicLiteratureReview2025} و پیشنهاد شده که معماری‌های ترکیبی که زنجیره‌ی بلوکی، رایانش ابری و رایانش مه و لبه را یکپارچه می‌کنند، و توزیع وظایف مرجع قابل اعتماد، نقطه شکست متمرکز را کاهش دهند \cite{rahmawatiagustinaSystematicLiteratureReview2025}. دوم، تحلیل‌های رسمی حریم خصوصی عمدتاً غیررسمی‌اند: از ابزار \lr{ProVerif}\LTRfootnote{\lr{ProVerif}} تنها در دو مطالعه استفاده شده و ابزارهایی مانند \lr{ROM}\LTRfootnote{\lr{ROM}}، \lr{ECDLP}\LTRfootnote{\lr{ECDLP}}، \lr{ECDH}\LTRfootnote{\lr{ECDH}}، منطق \lr{BAN}\LTRfootnote{\lr{BAN logic}}، \lr{AVISPA}\LTRfootnote{\lr{AVISPA}} و \lr{SPAN}\LTRfootnote{\lr{SPAN}} برای تحلیل حریم خصوصی طراحی نشده‌اند \cite{rahmawatiagustinaSystematicLiteratureReview2025}. سوم، مقیاس‌پذیری لغو کمتر بررسی شده است و مشخصه‌ی لغو \lr{W3C} هنوز استاندارد نیست \cite{mazzocca2025survey}؛ یک واحد کنارجاده‌ای می‌تواند فهرست لغو اعتبارنامه‌ها (\lr{CRL}\LTRfootnote{\lr{Certificate Revocation List (CRL)}}) را در منطقه‌ی شهری با کارایی پخش کند اما قابلیت اطمینان و ترافیک بیش از حد آن حل نشده است \cite{farsimadanReviewSecurityChallenges2025}. چهارم، تهدید کوانتومی موضوعی بلندمدت است: توابع درهم‌ساز کنونی ممکن است برای استفاده‌ی بلندمدت مناسب نباشند \cite{rahmawatiagustinaSystematicLiteratureReview2025}. پنجم، شفافیت ارزیابی کم است: بیشتر پژوهشگران مجموعه‌داده‌های ارزیابی خود را عمومی نمی‌کنند و بیشتر روش‌ها با شبیه‌سازهای ضعیف ارزیابی شده‌اند \cite{farsimadanReviewSecurityChallenges2025}، در حالی که سامانه‌ی \lr{DIVA} مجموعه‌داده‌ی سازگار با \lr{ETSI} منتشر کرده است \cite{feraudoDIVADIDbasedReputation2024} و آزمون‌های عملی بیشتری لازم است \cite{luPseudonymChangingSocial2012}. ششم، یادگیری فدرال\LTRfootnote{\lr{Federated Learning}} همچنان به ساختار سرور و مشتری وابسته است و در برابر مسموم‌سازی داده آسیب‌پذیر است \cite{farsimadanReviewSecurityChallenges2025} و هم‌پذیری چند ذی‌نفعی با دفتر کل توزیع‌شده می‌تواند پیچیدگی زمانی بیشتر ایجاد کند \cite{farsimadanReviewSecurityChallenges2025}. هفتم، حکمرانی میان مالکان خودروها و نهادهای دولتی دشوار است \cite{razaBlockchainBasedReputationTrust2024}. در ادامه، موضوعات آینده شامل ارتباطات \lr{V2X} نسل ششم، رایانش لبه، یادگیری فدرال و رایانش کوانتومی است \cite{hozouriOverviewVANETVehicular2023} و در تنگل نیز اجماع احتمالاتی سریع روی تنگل، اوراکل‌ها\LTRfootnote{\lr{Oracle}} و شاردینگ از پیشرفت‌های آینده‌اند \cite{sealeyIOTATangle202022}. در نهایت، یافته‌های این گزارش نشان می‌دهد که امنیت شبکه‌ی اقتضایی خودرویی مسئله‌ای چندبعدی است: از آسیب‌پذیری‌های ساختاری گرفته تا چالش‌های حکمرانی میان مالکان خودروها و نهادهای دولتی \cite{razaBlockchainBasedReputationTrust2024} و الزامات حقوقی مانند مقررات حفاظت از داده‌ها \cite{mazzocca2025survey}، و پیشرفت در آن مستلزم پژوهشی است که هم‌زمان ابعاد فنی و حقوقی را در بر بگیرد.
